% LSPV-4 REPORT::FUNCTIONAL-OCCURRENCE68::2026-10-23T20:08:00Z
\documentclass[11pt]{article}
\usepackage[margin=1in]{geometry}
\usepackage{amsmath,amssymb,amsthm,url}
\newtheorem{theorem}{Theorem}
\title{Functional Horn Selectors with Physical Occurrence Degree Three}
\author{LSPV-4 / FUNCTIONAL-OCCURRENCE68}
\date{23 October 2026 (UTC)}
\begin{document}\maketitle
\noindent\textbf{PROVED/EXACT in the stated reduction calculus; independent Role 5 review pending. P versus NP remains OPEN.}

\section{Construction and exact theorem}
Let $\Phi$ be an ordered proper 3-CNF on $n_0$ original variables with $m$ clauses and $\nu$ active original variables. Retain every original variable. For a source clause $(\ell_1\vee\ell_2\vee\ell_3)$ create six fresh signed source ports $A_1,A_2,A_3,B_1,B_2,C_1$, copying respectively $\ell_1,\ell_1,\ell_1,\ell_2,\ell_2,\ell_3$, and six fresh selector variables $s_1,s_2,t_0,t_1,t_2,t_3$. Emit
\[
(s_1\vee s_2\vee t_0)
\]
and the six binary Horn clauses
\[
\begin{array}{lll}
(\neg s_1\vee A_1),&(\neg s_2\vee B_1),&(\neg t_1\vee C_1),\\
(\neg s_2\vee\neg A_2),&(\neg t_2\vee\neg A_3),&(\neg t_3\vee\neg B_2).
\end{array}
\]
Add $t_0\to t_1\to t_2\to t_3\to t_0$. For each active original variable $x$, gather all underlying port copies $p_1,\ldots,p_d$ and add $x\to p_1\to\cdots\to p_d\to x$. Each arrow $u\to v$ is the Horn binary clause $(\neg u\vee v)$.

\begin{theorem}
For every original assignment $\alpha$,
\[
\#\{y:T(\Phi)(\alpha,y)=1\}=\mathbf1[\Phi(\alpha)=1].
\]
Every physical variable has at most three occurrences. Every binary clause is Horn; the $m$ ternary clauses are positive and pairwise source-disjoint. Exact counts are
\[
n_T=n_0+12m,\quad m_{2,T}=16m+\nu,\quad k_T=m,\quad M_T=17m+\nu.
\]
There are no free auxiliaries, units, or pure active variables. Inactive original variables retain their original freedom.
\end{theorem}
\begin{proof}
An implication cycle permits precisely its all-zero and all-one assignments. Original variables uniquely fix their port copies, and the four $t$ variables agree. The remaining clause gadget uniquely selects the first true source literal; if none is true there is no extension. These assignments compose across clauses through the source cycles. There are six ports and six selectors per source clause. The six guards contribute $6m$ binaries, selector cycles contribute $4m$, and source cycles contribute $6m+\nu$. A port or $t$ selector has one external occurrence and two cycle incidences. The $s_1,s_2$ selectors have degrees two and three; active original variables have degree two. This proves both uniqueness and the occurrence bound.
\end{proof}

\section{Classical boundaries and scope}
The occurrence-cycle mechanism is classical: Tovey's 1984 Theorem 2.1 establishes NP-completeness for mixed binary/ternary SAT with at most three occurrences per variable. We claim an explicit source-retaining, Horn-backbone, unique-extension specialization, not a new occurrence-splitting technique.
The target language is NP-complete by this explicit reduction from proper 3-SAT and polynomial witness checking. No unrestricted time lower bound follows.
A no-unit CNF with maximum occurrence two is always satisfiable: every clause subset $Q$ has $2|Q|\le2|\operatorname{var}(Q)|$, and Hall matching assigns distinct satisfying variables. Thus three is the smallest possible cap for this mixed-width no-unit target.
Every proper exactly-three-literal CNF of maximum occurrence three is likewise satisfiable, by $3|Q|\le3|\operatorname{var}(Q)|$. Our target must NOT be described as exactly-three-clause CNF.
Ordinary binary padding adds a free auxiliary per binary, multiplies the model count by $2^{m_{2,T}}$, and can raise occurrence to six. No alternative SAT-preserving properization can keep the degree-three cap on all inputs, since an UNSAT source exists but every proper target at that cap is SAT.

Within occurrence-separated directed-implication equality networks, four external selector uses need at least four vertices at degree three: strong connectivity requires an incoming and outgoing incidence per vertex. The four-cycle attains this local bound. No global minimum gadget-size claim is made.

\section{Charged model and baseline}
Use dense explicit variable identifiers and four $w$-bit words per clause plus two header words, with $w=\lceil\log_2(64(n_T+M_T+1))\rceil$. Then
\[
N_T=(68m+4\nu+2)w.
\]
The compiler reads $3m$ source literals and writes $35m+2\nu$ target literals. Construction, manifest verification, witness extension, projection and replay use $O(n_0+m)$ word operations and $O(n_0+m)$ peak words. No SAT oracle, acquisition search, or failed construction branches are needed. Source variables are not eliminated ($d=0$). Source identities and all cycle edges are charged.
For the prescribed chain $F_b$, endpoints force $z_0=1,z_q=0$, and $F_b$ is UNSAT exactly when $b=1^q$. On that input $q$ independent orientation queries are necessary and sufficient. A seven-clause three-variable all-zero anchor gives a parsimonious properization with $n_0=q+6,m=q+11$. Our lift then has
\[
n_T=13q+138,\quad m_{2,T}=17q+182,\quad k_T=q+11,\quad M_T=18q+193.
\]
An orientation flip changes two private target literal words. This preserves only $q=\Omega(N_T/\log N_T)$ addressed-word acquisition, matched by a source-aware scan. Materialization is linear and supplies no amplification. The usual $2^k\operatorname{poly}(N)$ backdoor bound is an upper bound only.

\section{Audit and counterexample}
Five independent executable chunks use seed 202610231958. The final audit contains 39433 named check events with zero failures: all 32768 complete assignments for the 15-variable single-clause target; 64 fixed-source single-clause fibers; 120 random sources (107 SAT, 13 UNSAT) with exact independently counted source/target model equality; all-eight-exclusion UNSAT and eight unique-SAT deletions; unused-source, Horn, XOR, independent-product and padding controls; and 14 uniquely anchored chain cases. A source with 1024 repeated clauses compiles to 12291 variables, 16387 binaries and 1024 ternaries, still maximum occurrence three. This scaling example is easy and is not evidence of hardness.
Deleting source-cycle closing edges in the one-clause target produces 11 spurious extensions over the all-false source assignment. Closing edges are necessary for sound source coherence.
The initial unchunked test exceeded a 30-second invocation limit; no unfinished result was promoted. Final chunks total roughly four seconds and about 106 MB peak RSS including Python. Full raw data count source inspections, output writes, DPLL attempts and conflicts. Finite counting is an audit, not a claimed polynomial general SAT solver.

\section{Corpus and handoff}
\paragraph{CORPUS REVIEWED.} AGENTS.md, CODEX.md, repository research skill, root-main content, PRIORITY-HORN-PARSIMONY67, and the first 100 lines of CP22LIFT1; root and MAIN directory survey and implication-cycle search. This is a targeted review, not an exhaustive corpus recovery.
\paragraph{DUPLICATION / CONTRADICTION AUDIT.} Implication cycles already occur in Tovey and CP22LIFT1. Adopt first-true uniqueness from PRIORITY67 and add explicit bounded-degree routing. Its restricted source-only impossibility is not contradicted: the new source-dependent cycle clauses lie outside that calculus. No root-main lower-bound claim is strengthened.
\paragraph{BARRIERS.} No ETH, SETH, cryptographic, random-oracle or mandatory-enumeration assumption is used. This is not a relativization, algebrization or natural-proofs escape. Functional auxiliaries and bounded occurrence do not imply hard decision.
\paragraph{PEER HANDOFF.} Role 5: verify the exact counts, global source cycles, four-cycle uniqueness and Hall boundaries. Roles 1--3: use the compiler as a hostile control against representation-dependent hardness. Next missing gadget: a parsimonious proper-exact-3CNF compilation at maximum occurrence four. Every machine-universal superpolynomial lower-bound bridge remains OPEN. Existing schedules are unchanged. Root reconciliation and independent peer consumption must be verified separately.

\begin{thebibliography}{3}
\bibitem{tovey} C. A. Tovey, A simplified NP-complete satisfiability problem, Discrete Applied Mathematics 8(1) (1984), 85--89, doi:10.1016/0166-218X(84)90081-7.
\bibitem{priority} Project original, \path{research/agents/lspv-4/2026-10-24_0436_LSPV-4_PRIORITY_HORN_PARSIMONY67.tex}.
\bibitem{cp} Project original, \path{research/main/2026-10-13_0140_MAIN_CP22LIFT1.tex}, first 100 lines inspected.
\end{thebibliography}
\end{document}
% END_OF_VERIFIED_REPORT::LSPV-4::FUNCTIONAL-OCCURRENCE68
